\documentclass{article}

\usepackage{amsmath,amssymb}
\usepackage{xurl}
\usepackage[hidelinks]{hyperref}

% Shared commands for notes in docs/paper-gaps/.

\DeclareUnicodeCharacter{2080}{\ifmmode{}_{0}\else\textsubscript{0}\fi}
\DeclareUnicodeCharacter{2081}{\ifmmode{}_{1}\else\textsubscript{1}\fi}
\DeclareUnicodeCharacter{2082}{\ifmmode{}_{2}\else\textsubscript{2}\fi}
\DeclareUnicodeCharacter{2099}{\ifmmode{}_{n}\else\textsubscript{n}\fi}

\newcommand{\Lean}{\textsc{Lean}}
\ExplSyntaxOn
\NewDocumentCommand{\leanid}{m}{
  \ifmmode
    \textsf{\detokenize{#1}}
  \else
    \group_begin:
    \urlstyle{sf}
    \tl_set:Nn \l_tmpa_tl {#1}
    \tl_replace_all:Nnn \l_tmpa_tl {\_} {_}
    \exp_args:NV \path \l_tmpa_tl
    \group_end:
  \fi
}
\ExplSyntaxOff
\providecommand{\tr}{\operatorname{tr}}
\newcommand{\tnleanIssueBase}{https://github.com/LionSR/TNLean/issues/}
\newcommand{\tnleanPrBase}{https://github.com/LionSR/TNLean/pull/}
\newcommand{\tnleanDocBase}{https://sirui-lu.com/TNLean/docs/}
\newcommand{\ghissue}[1]{\href{\tnleanIssueBase#1}{GitHub issue \##1}}
\newcommand{\ghpr}[1]{\href{\tnleanPrBase#1}{PR \##1}}
\newcommand{\doclink}[1]{\href{\tnleanDocBase#1.html}{\path{#1}}}

% Dirac notation and the identity operator, as in blueprint/src/macros/common.tex.
% \providecommand keeps note-local definitions authoritative.
\usepackage{dsfont}
\providecommand{\ket}[1]{|#1\rangle}
\providecommand{\bra}[1]{\langle #1|}
\providecommand{\braket}[2]{\langle #1 | #2 \rangle}
\providecommand{\ketbra}[2]{|#1\rangle\!\langle #2|}
\providecommand{\Id}{\mathds{1}}
\providecommand{\one}{\mathds{1}}
\providecommand{\C}{\mathbb{C}}
\providecommand{\MN}[1]{M_{#1}(\C)}
\providecommand{\MD}{M_D(\C)}

% External referencing for paper sources.  The build helper writes sanitized
% auxiliary files under build/paper-aux/ and excludes duplicated source labels.
\usepackage{xr}

% Import a generated paper auxiliary file with a caller-supplied label prefix.
% The candidate paths support builds from docs/paper-gaps/, the repository root,
% and build/paper-aux/ itself.
\newcommand{\importpaperaux}[2]{%
  \IfFileExists{../../build/paper-aux/#2.aux}{%
    \externaldocument[#1]{../../build/paper-aux/#2}%
  }{%
    \IfFileExists{build/paper-aux/#2.aux}{%
      \externaldocument[#1]{build/paper-aux/#2}%
    }{%
      \IfFileExists{#2.aux}{%
        \externaldocument[#1]{#2}%
      }{}%
    }%
  }%
}

% General external reference macros
\newcommand{\paperref}[2]{\ref{#1-#2}}
\newcommand{\papereqref}[2]{(\ref{#1-#2})}

% CPSV16 source (1606.00608 / MPDO-22-12-17-2.tex) external document and shortcuts
\importpaperaux{cpsv16-}{1606.00608/MPDO-22-12-17-2-tnlean}
\newcommand{\cpsvref}[1]{\paperref{cpsv16}{#1}}
\newcommand{\cpsveqref}[1]{\papereqref{cpsv16}{#1}}

% Verdict marker: \gapnote{<kind>}{<status>}. Kind is the policy
% classification (clarification, local-correction, scope-restriction,
% unfaithful, false-source, open-gap); status is open, wip (elimination
% underway), resolved, or historical. Machine-read by the site index;
% prints a verdict line.
\newcommand{\gapnote}[2]{\par\noindent\textbf{Verdict:} #1 (#2).\par}


\title{The Fusion Gauge of the Stabilizer Reconstruction Formula}
\author{}
\date{2026-09-26}

\begin{document}
\maketitle
\gapnote{local-correction}{resolved}

\section{The source statement}

Let a group $G$ act transitively on a set $\mathsf X$ of blocks, fix
$x_0\in\mathsf X$ with stabilizer $H$, and choose representatives $k_x$
with $k_xx_0=x$ and $k_{x_0}=e$. Let $\omega$ be a three-cocycle on $G$.
Ref.~\cite[source lines~740--747]{GarreRubio2023MPOSym} shows that the
restriction $\omega|_H$ is a coboundary $d\beta$, and that in the fusion
gauge fixed by $\beta$ the basepoint restriction $\psi=L/\beta$ is a
two-cocycle on $H$. Equation~(20) of the same reference, at source
lines~752--759, then reconstructs every component of the $L$-symbol:
\begin{align}
    L^x_{g_1,g_2}
        &=\frac{\omega(g_1g_2k_x,h_2^{-1},h_1^{-1})\,\psi(h_1,h_2)}
          {\omega(g_1,g_2k_x,h_2^{-1})\,\omega(g_1,g_2,k_x)},
    \label{eq:glm23_eq20}
\end{align}
where $h_2=k_{g_2x}^{-1}g_2k_x$ and $h_1=k_{g_1g_2x}^{-1}g_1k_{g_2x}$.
Line~761 says that ``given $\omega$ and $\beta$ by fixing some gauge'' the
solutions of the compatibility equation
\begin{align}
    L^x_{g,hk}L^x_{h,k}
        &=\omega(g,h,k)L^{kx}_{g,h}L^x_{gh,k}
    \label{eq:glm23_compatibility}
\end{align}
form a torsor under $H^2(H,\mathbb C^\times)$.

\section{The gap}

Formula~(\ref{eq:glm23_eq20}) with a two-cocycle $\psi$ satisfies
(\ref{eq:glm23_compatibility}) when $\omega(a,b,c)=1$ for all
$a,b,c\in H$. It does not do so for a general three-cocycle with
$\omega|_H=d\beta$. At $x=x_0$ and $a,b\in H$ the formula gives
$L^{x_0}_{a,b}=F(a,b)\psi(a,b)$ with
\begin{align}
    F(a,b)
        &=\frac{\omega(ab,b^{-1},a^{-1})}{\omega(a,b,b^{-1})\,\omega(a,b,e)},
\end{align}
while (\ref{eq:glm23_compatibility}) restricted to $x_0$ and $H$ forces
$\omega|_H=d(L^{x_0}|_H)=dF$. This can fail.

\paragraph{Example.} Let $G=H=\mathbb Z_3=\{e,s,s^2\}$ act on a single
block, so $k=e$ and $h_1=g_1$, $h_2=g_2$. Let $\beta(s,s)=-1$ and
$\beta=1$ at every other pair, and put $\omega=d\beta$, with
$(d\beta)(g,h,k)=\beta(g,hk)\beta(h,k)/(\beta(g,h)\beta(gh,k))$. Then
$\omega(s,s,s^2)=-1$, while $F(s,s)=F(s^2,s^2)=-1$ and $F=1$ at the other
pairs, so $dF=1$ identically. Hence for every two-cocycle $\psi$ the
formula violates (\ref{eq:glm23_compatibility}) at $(g,h,k)=(s,s,s^2)$.
Numerical checks on $A_4\supset V_4$ and $S_3\supset A_3$, with $\omega$
twisted by a random normalized coboundary, show violations of order one in
the same way.

\section{The local fix}

Before applying (\ref{eq:glm23_eq20}), change the fusion gauge so that
$\omega$ is one on $H\times H\times H$. Such a gauge always exists: if
$\omega|_H=d\beta_0$ with $\beta_0$ on $H$, extend $\beta_0^{-1}$ to $G$ by
one outside $H\times H$ and apply the resulting fusion gauge. In that gauge
$F=1$ on $H$, the basepoint restriction of (\ref{eq:glm23_eq20}) is $\psi$
itself, and a direct four-term three-cocycle computation shows that
(\ref{eq:glm23_eq20}) satisfies (\ref{eq:glm23_compatibility}); no
normalization of $\omega$ is needed. Transporting back by the inverse
fusion gauge gives a solution for the original $\omega$. Equivalently, one
may keep $\omega$ and replace $\psi$ in (\ref{eq:glm23_eq20}) by
$\psi\beta F^{-1}$, which matches the source's reading $L|_H=\psi\beta$.

The torsor statement of line~761 is correct for classes under
action-tensor gauges with the fusion gauge held fixed, as the source says.
If fusion gauges preserving $\omega$ are also allowed, a two-cocycle $\beta$
on $G$ acts, and the classes collapse to the quotient of
$H^2(H,\mathbb C^\times)$ by the restriction of
$H^2(G,\mathbb C^\times)$; this remark is not formalized.\footnote{Formal declarations:
\leanid{TNLean.Algebra.StabilizerRepresentatives.reconstructedLSymbol},
\leanid{TNLean.Algebra.StabilizerRepresentatives.reconstructedLSymbol\_isCompatible},
\leanid{TNLean.Algebra.StabilizerRepresentatives.exists\_actionGauge\_eq\_reconstructedLSymbol},
\leanid{TNLean.Algebra.StabilizerRepresentatives.h2EquivSolutionClasses} and
\leanid{TNLean.Algebra.ScalarThreeCochain.exists\_fusionGauge\_eq\_one\_of\_isTrivialGaugeClass\_comap}
in \doclink{TNLean/Algebra/StabilizerCocycleReconstruction}.}

\textbf{Verdict: resolved local fix.} Formula~(\ref{eq:glm23_eq20}) is
correct in the fusion gauge in which $\omega$ is trivial on the stabilizer,
and that gauge always exists. The classification and the torsor statement
hold as stated for action-tensor gauge classes.

\bibliographystyle{alpha}
\bibliography{references}

\end{document}
